\documentclass[11pt,a4paper]{article}

\usepackage[T1]{fontenc}
\usepackage[utf8]{inputenc}
\IfFileExists{lmodern.sty}{\usepackage{lmodern}\usepackage{microtype}}{}
\IfFileExists{ngerman.ldf}{\usepackage[ngerman]{babel}}{}
\usepackage[top=1.9cm,bottom=1.7cm,left=2.3cm,right=2.3cm]{geometry}

\setlength{\parskip}{0.55em}
\setlength{\parindent}{0pt}

\begin{document}
\pagestyle{empty}

%% ============================================================
%% Motivationsschreiben [First] [Last]
%% WSL, Landschaftsdynamik / Landnutzungssysteme, DISCLOSE
%% Gestrafft: Motivation und Passung, vier Absaetze.
%% Ohne Gedankenstriche und ohne Fettdruck.
%% ============================================================

\begin{flushleft}
[First] [Last]\\
Zürich\\
{[Phone]}\\
{[your.email@example.com]}
\end{flushleft}

\begin{flushright}
Eidg. Forschungsanstalt WSL\\
Beatrice Lamprecht, Human Resources\\
Zürcherstrasse 111\\
8903 Birmensdorf
\end{flushright}

Zürich, 28. August 2026

\vspace{0.2em}

Bewerbung als Wissenschaftliche Assistenz Sozialwissenschaftliche Forschung, DISCLOSE

Sehr geehrte Frau Lamprecht, sehr geehrte Frau Hersperger

Drei Jahre lang habe ich an der ETH Zürich in der Gruppe Raumentwicklung und Stadtpolitik von Prof. Dr. David Kaufmann untersucht, wie Städte Nachhaltigkeitsziele formulieren und wessen Anliegen sich dabei durchsetzen. Rund 300 regionale Planungs- und Politikdokumente habe ich dafür codiert und daraus eine Typologie städtischer Nachhaltigkeitspolitik entwickelt; aus derselben Arbeit ist eine Ko-Autorenschaft in \textit{Nature Sustainability} entstanden, die zeigt, wie weit versprochene und tatsächliche Beteiligung auseinanderfallen. Wer in einem Planungsprozess vorkommt und wer nicht, beschäftigt mich seitdem. Genau diese Frage stellt DISCLOSE unter neuen Bedingungen, und deshalb bewerbe ich mich.

Die Ausschreibung trifft dabei dreifach zu. Inhaltlich sind städtische Räume, Raumplanung und Bodennutzung seit dem Studium mein Gebiet. Methodisch ist die selbstständige Auswertung qualitativer Daten aus Interviews und Dokumenten meine tägliche Arbeit, codiert in MAXQDA, an der ETH ebenso wie bei Interface, Sotomo und der LMU. Und das Interesse ist kein neues: Wie Regeln, Verfahren und neuerdings Datenmodelle bestimmen, was in der Planung sichtbar wird, zieht sich durch alles, was ich bisher gemacht habe.

Bei Interface Politikstudien habe ich Evaluationen für Bund, Kantone und Gemeinden erarbeitet, mit halbstrukturierten Interviews, Fokusgruppen und Inhaltsanalysen von Verwaltungsunterlagen, und dabei untersucht, wie dieselbe Vorgabe je nach Stelle unterschiedlich umgesetzt wird. Das ist im Kern die Frage nach dem fachlichen Ermessen, und sie interessiert mich an DISCLOSE am stärksten: Sobald ein Verfahren digital abgebildet wird, wandert dieser Spielraum in Feldtypen, Pflichtangaben und Validierungsregeln, die sich später kaum revidieren lassen. Dass Ihr Projekt Digitalisierung deshalb als soziotechnische Reorganisation der Planungspraxis fasst und nicht als Einführung neuer Werkzeuge, halte ich für den richtigen Ausgangspunkt. Bei Sotomo habe ich denselben Zusammenhang im kommunalen Massstab untersucht: wer die digitalen Dienstleistungen der Stadt Zürich nutzt, wer nicht, und woran es liegt.

Beobachtet habe ich bisher im Feld, nicht in Planungsbüros; die organisationsethnografische Arbeit wäre für mich eine Erweiterung, auf die ich mich freue. Die Stelle ist als Assistenz ausgeschrieben, und genau als solche bewerbe ich mich: Interviewmaterial aufzubereiten, Literatur zu sichten, Workshops vorzubereiten und die Projektorganisation mitzutragen gehört für mich zum Beitrag ebenso wie die eigene Analyse.

Zur Einordnung: Mein ETH-Master heisst Comparative and International Studies, inhaltlich auf empirische Sozialforschung und Stadtforschung ausgerichtet. Meine derzeitige Stelle an der LMU München, an der ich qualitative und quantitative Sozialforschung zu staatlichen Massnahmen betreibe, ist befristet und war von Anfang an als Übergang angelegt; ein Pensum von 60 bis 80 Prozent ab Herbst passt daher gut. Deutsch ist meine Arbeitssprache, Englisch meine Muttersprache. Über ein Gespräch freue ich mich sehr.

\vspace{0.3em}

Freundliche Grüsse

\vspace{0.8em}

[First] [Last]

\end{document}
